\begin{conf-abstract}
{A first look at the dissolved gas composition of offshore freshened groundwater (IODP 501, New England continental shelf, USA)}
{David V. Bekaert\textsuperscript{1}, Alizé Longeau\textsuperscript{1}, Paul Moser Roeggla\textsuperscript{2}, Rolf Kipfer\textsuperscript{2}, Alan Seltzer\textsuperscript{3,4}, Israela Musan\textsuperscript{4}, Douglas Kip Solomon\textsuperscript{5}, William Davis Mace\textsuperscript{5}}
{(1) CRPG-Lorraine University, Nancy (France)\\
(2) Eawag, Swiss Federal Institute of Aquatic Science and Technology\\
(3) University College Dublin, Ireland\\
(4) Woods Hole Oceanographic Institution, MA, USA\\
(5) The University of Utah, UT, USA
}
{Freshened porewater in marine sediments (characterized by its lower salinity relative to the overlying seawater) has been recognized since the 1960s and is referred to as offshore freshened groundwater (OFG) [1–2]. Evidence from scientific ocean drilling and hydrological modeling suggests that 105–106 km3 of OFG may be stored within continental shelf aquifers worldwide [3]. The possibility that OFG constitutes a substantial, yet largely untapped, source of freshwater has stimulated growing scientific interest over the past decade, with important implications for coastal and marine hydrogeology, paleoclimate and paleoenvironmental reconstructions, and groundwater-dependent coastal ecosystems. However, the origin(s), emplacement history, and dynamics of these shoreline-crossing groundwater systems remain poorly constrained.

Salinity data from the Atlantic continental shelf off New England have indicated that the present-day freshwater/saltwater interface is far out of equilibrium with modern sea-level conditions [4]. To shed light on the origin and evolution history of this hydrological system, the IODP Expedition 501 ``New England Shelf Hydrogeology'' cored and sampled the subseafloor offshore Massachusetts, USA (\url{https://www.ecord.org/expedition501/}), from May to August 2025. For the first time, they conducted a series of successful pumping tests on OFG bodies located up to ~400 m below the seafloor. To monitor these pumping tests, we used a miniRuedi to measure dissolved gas concentrations in real time, enabling (i) assessment of the representativeness of samples collected by the multiple scientific teams, (ii) evaluation of the hydraulic confinement of the targeted aquifers, and (iii) qualitative estimation of groundwater residence times. The fun started right with the very first pumping test, when we realized that the best we were going to get was brackish water packed with fine particles.

This presentation will walk you through the motivation, strategy, and challenges of this project. We will share some of the general results and preliminary interpretations covering both the technical side of the pumping tests and the preliminary scientific insights they have revealed about the aquifers we explored.

[1] Kohout (1964), The flow of fresh water and salt water in the Biscayne aquifer of the Miami area, FloridaRep., USGS Water Supply, Washington, DC.\newline
[2] Micallef, Person, Berndt, Bertoni, Cohen, et al. (2021), Offshore freshened groundwater in Continental margins, Reviews of Geophysics, 59(1), doi:10.1029/2020rg000706.\newline
[3] Zamrsky, Essink, Sutanudjaja, Van Beek, and Bierkens (2021), Offshore fresh groundwater in coastal unconsolidated sediment systems as a potential fresh water source in the 21st century, Environmental Research Letters, doi:10.1088/1748-9326/ac4073.\newline
[4] Person, M., Dugan, B., Swenson, J. B., Urbano, L., Stott, C., Taylor, J., \& Willett, M. (2003). Pleistocene hydrogeology of the Atlantic continental shelf, New England. Geological Society of America Bulletin, 115(11), 1324-1343.
}
\end{conf-abstract}
